\documentclass[12pt,a4paper,openany,oneside]{book}

% --- Paquetes ---
\usepackage[utf8]{inputenc}
\usepackage[spanish, es-noshorthands]{babel}
\usepackage{amsmath, amsfonts, amssymb}
\usepackage{graphicx}
\usepackage{geometry}
\usepackage{hyperref}
\usepackage{booktabs}
\usepackage{fancyhdr}
\usepackage{listings}
\usepackage{xcolor}
\usepackage{tikz}
\usepackage{pgfplots}
\usetikzlibrary{arrows.meta, positioning, shapes.geometric}
\pgfplotsset{compat=1.18}

\geometry{margin=2.5cm}

% --- Colores y estilos para código Python ---
\definecolor{codegreen}{rgb}{0,0.6,0}
\definecolor{codegray}{rgb}{0.5,0.5,0.5}
\definecolor{codepurple}{rgb}{0.58,0,0.82}
\definecolor{backcolour}{rgb}{0.95,0.95,0.92}

\lstset{
    language=Python,
    backgroundcolor=\color{backcolour},   
    commentstyle=\color{codegreen},
    keywordstyle=\color{blue},
    numberstyle=\tiny\color{codegray},
    stringstyle=\color{codepurple},
    basicstyle=\ttfamily\small,
    breaklines=true,
    frame=single,
    showstringspaces=false
}

% --- Estilo de cabecera y pie ---
\pagestyle{fancy}
\fancyhf{}
\renewcommand{\headrulewidth}{0.4pt}
\renewcommand{\footrulewidth}{0.4pt}
\fancyhead[L]{\small Carlos César Sánchez Coronel}
\fancyhead[R]{\small Sesión 12: Autovalores y Autovectores}
\fancyfoot[L]{\small Carlos César Sánchez Coronel}
\fancyfoot[C]{\thepage}

% --- Portada ---
\title{
    \textbf{Sesión 12} \\ 
    \Huge \textbf{AUTOVALORES Y AUTOVECTORES} \\
    \large Las direcciones clave de las transformaciones lineales \\
    \normalsize Aplicaciones en PCA, PageRank y análisis de redes neuronales
}
\author{
    \small Por \\ \vspace{0.2cm}
    \Large \textsc{Carlos César Sánchez Coronel}
}
\date{2026}

\begin{document}

\maketitle
\tableofcontents
\addcontentsline{toc}{chapter}{Índice general}

% ------------------------------------------------------------------
% CAPÍTULO 12: SESIÓN 12 - AUTOVALORES Y AUTOVECTORES
% ------------------------------------------------------------------
\chapter{Autovalores y Autovectores: Las Direcciones Clave}

Los autovalores y autovectores (o valores y vectores propios) son uno de los conceptos más poderosos y elegantes del álgebra lineal. Nos permiten ``descomponer'' una transformación lineal en sus direcciones fundamentales: aquellas que, al aplicar la transformación, solo se estiran o encogen, pero no cambian de dirección. Esta idea tiene aplicaciones profundas en inteligencia artificial: desde el análisis de componentes principales (PCA) para reducir dimensionalidad, hasta el algoritmo PageRank de Google y el estudio de la dinámica de aprendizaje en redes neuronales.

\section{Definición y significado geométrico}

\subsection{Definición formal}
Sea $\mathbf{A} \in \mathbb{R}^{n \times n}$ una matriz cuadrada. Un escalar $\lambda \in \mathbb{C}$ (generalmente real si la matriz es simétrica) y un vector $\mathbf{v} \in \mathbb{C}^n$ no nulo que satisfacen:
\[
\mathbf{A} \mathbf{v} = \lambda \mathbf{v}
\]
se denominan, respectivamente, \textbf{autovalor} y \textbf{autovector} de $\mathbf{A}$.

En otras palabras, cuando aplicamos la transformación lineal representada por $\mathbf{A}$ al vector $\mathbf{v}$, el resultado es simplemente el mismo vector multiplicado por el escalar $\lambda$. Por tanto, la dirección de $\mathbf{v}$ permanece inalterada (o se invierte si $\lambda$ es negativo), y la magnitud se escala por $|\lambda|$.

\subsection{Interpretación geométrica}
Para visualizarlo, consideremos una matriz $\mathbf{A}$ en $\mathbb{R}^2$. Un autovector $\mathbf{v}$ señala una dirección que no se rota bajo la transformación, solo se estira. Por ejemplo, si $\mathbf{A}$ es una matriz de rotación pura (excepto rotaciones de 0° o 180°), no tiene autovectores reales porque todas las direcciones cambian. En cambio, una matriz de escalado tiene autovectores en las direcciones de los ejes.

\begin{center}
\begin{tikzpicture}
\draw[->] (-1,0) -- (3,0) node[right] {$x$};
\draw[->] (0,-1) -- (0,3) node[above] {$y$};
\draw[->, thick, blue] (0,0) -- (1,2) node[midway, above left] {$\mathbf{v}$};
\draw[->, thick, red] (0,0) -- (2,1) node[midway, below right] {$\mathbf{w}$};
\node at (0.5,1) {$\bullet$};
\end{tikzpicture}
\captionof{figure}{Si $\mathbf{v}$ es autovector, $\mathbf{A}\mathbf{v}$ es paralelo a $\mathbf{v}$.}
\end{center}

\section{Cálculo de autovalores y autovectores}

\subsection{Ecuación característica}
De la definición $\mathbf{A}\mathbf{v} = \lambda \mathbf{v}$, reordenamos:
\[
\mathbf{A}\mathbf{v} - \lambda \mathbf{v} = \mathbf{0} \quad \Rightarrow \quad (\mathbf{A} - \lambda \mathbf{I}) \mathbf{v} = \mathbf{0}
\]
Para que exista un vector $\mathbf{v} \neq \mathbf{0}$, la matriz $\mathbf{A} - \lambda \mathbf{I}$ debe ser singular, es decir, su determinante debe ser cero:
\[
\det(\mathbf{A} - \lambda \mathbf{I}) = 0
\]
Esta ecuación se llama \textbf{ecuación característica} y es un polinomio de grado $n$ en $\lambda$, cuyas raíces son los autovalores.

\subsection{Ejemplo manual para una matriz $2\times 2$}
Sea
\[
\mathbf{A} = \begin{pmatrix} 4 & 1 \\ 2 & 3 \end{pmatrix}
\]
Calculamos $\mathbf{A} - \lambda \mathbf{I} = \begin{pmatrix} 4-\lambda & 1 \\ 2 & 3-\lambda \end{pmatrix}$.
Su determinante:
\[
(4-\lambda)(3-\lambda) - 2 = \lambda^2 - 7\lambda + 12 - 2 = \lambda^2 - 7\lambda + 10 = (\lambda - 5)(\lambda - 2)
\]
Por tanto, los autovalores son $\lambda_1 = 5$ y $\lambda_2 = 2$.

Para $\lambda_1 = 5$, resolvemos $(\mathbf{A} - 5\mathbf{I})\mathbf{v} = \mathbf{0}$:
\[
\begin{pmatrix} -1 & 1 \\ 2 & -2 \end{pmatrix} \begin{pmatrix} x \\ y \end{pmatrix} = \begin{pmatrix} 0 \\ 0 \end{pmatrix}
\]
La primera ecuación: $-x + y = 0 \Rightarrow y = x$. Así, un autovector es $\mathbf{v}_1 = \begin{pmatrix} 1 \\ 1 \end{pmatrix}$ (cualquier múltiplo).

Para $\lambda_2 = 2$:
\[
\begin{pmatrix} 2 & 1 \\ 2 & 1 \end{pmatrix} \begin{pmatrix} x \\ y \end{pmatrix} = \begin{pmatrix} 0 \\ 0 \end{pmatrix}
\]
La primera ecuación: $2x + y = 0 \Rightarrow y = -2x$. Un autovector es $\mathbf{v}_2 = \begin{pmatrix} 1 \\ -2 \end{pmatrix}$.

Verificamos: $\mathbf{A}\mathbf{v}_1 = \begin{pmatrix} 4 & 1 \\ 2 & 3 \end{pmatrix} \begin{pmatrix} 1 \\ 1 \end{pmatrix} = \begin{pmatrix} 5 \\ 5 \end{pmatrix} = 5 \mathbf{v}_1$.

\subsection{Propiedades importantes}
\begin{itemize}
    \item La \textbf{traza} de $\mathbf{A}$ (suma de la diagonal) es igual a la suma de los autovalores (contando multiplicidades): $\operatorname{tr}(\mathbf{A}) = \sum_i \lambda_i$.
    \item El \textbf{determinante} de $\mathbf{A}$ es el producto de los autovalores: $\det(\mathbf{A}) = \prod_i \lambda_i$.
    \item Si $\mathbf{A}$ es simétrica ($\mathbf{A}^\top = \mathbf{A}$), todos los autovalores son reales y los autovectores pueden elegirse ortonormales.
    \item Si $\mathbf{A}$ es definida positiva (o semidefinida positiva), todos sus autovalores son positivos (o no negativos).
\end{itemize}

\section{Diagonalización}

Si $\mathbf{A} \in \mathbb{R}^{n \times n}$ tiene $n$ autovectores linealmente independientes, podemos formar la matriz $\mathbf{P}$ cuyas columnas son esos autovectores, y entonces:
\[
\mathbf{A} = \mathbf{P} \mathbf{\Lambda} \mathbf{P}^{-1}
\]
donde $\mathbf{\Lambda}$ es una matriz diagonal con los autovalores en la diagonal. Esto se llama \textbf{diagonalización} y simplifica enormemente el cálculo de potencias de $\mathbf{A}$:
\[
\mathbf{A}^k = \mathbf{P} \mathbf{\Lambda}^k \mathbf{P}^{-1}
\]
En particular, si $\mathbf{A}$ es simétrica, $\mathbf{P}$ puede ser ortogonal ($\mathbf{P}^{-1} = \mathbf{P}^\top$).

\section{Aplicaciones en Inteligencia Artificial}

\subsection{PCA (Análisis de Componentes Principales)}
El PCA es una técnica de reducción de dimensionalidad que busca las direcciones de máxima varianza en los datos. Dado un conjunto de datos centrado $\mathbf{X} \in \mathbb{R}^{n \times d}$ (cada fila es una muestra), la matriz de covarianza muestral es:
\[
\boldsymbol{\Sigma} = \frac{1}{n-1} \mathbf{X}^\top \mathbf{X}
\]
Los autovalores de $\boldsymbol{\Sigma}$ representan la varianza a lo largo de los autovectores correspondientes. Los autovectores con mayores autovalores son las \textbf{componentes principales}. Proyectando los datos sobre esos autovectores, reducimos la dimensión reteniendo la mayor parte de la varianza.

Pasos del PCA:
\begin{enumerate}
    \item Centrar los datos (restar la media).
    \item Calcular la matriz de covarianza.
    \item Obtener autovalores y autovectores (por ejemplo, con SVD o \texttt{eig}).
    \item Ordenar autovalores de mayor a menor y seleccionar los $k$ primeros autovectores.
    \item Proyectar los datos: $\mathbf{X}_{\text{proj}} = \mathbf{X} \mathbf{V}_k$, donde $\mathbf{V}_k$ contiene los $k$ autovectores seleccionados.
\end{enumerate}

\subsection{PageRank de Google}
El algoritmo PageRank modela la importancia de las páginas web mediante un grafo dirigido. La matriz de adyacencia se normaliza para obtener una matriz estocástica $\mathbf{G}$ (con columnas que suman 1). El vector de ranking $\mathbf{r}$ es el autovector principal de $\mathbf{G}$, es decir, satisface:
\[
\mathbf{G} \mathbf{r} = \mathbf{r}
\]
(o con un factor de amortiguamiento). La solución es el autovector asociado al autovalor $\lambda = 1$. Se calcula mediante el método de la potencia (iteración).

\subsection{Análisis de redes neuronales}
En deep learning, los autovalores de la matriz de pesos pueden dar información sobre la dinámica del aprendizaje. Por ejemplo, en redes recurrentes (RNN), los autovalores de la matriz de pesos recurrente determinan la estabilidad y la capacidad de retener memoria a largo plazo (gradientes explosivos o desvanecientes). También se usa el análisis de autovalores del Hessiano para entender la curvatura de la función de pérdida.

\section{Python: Cálculo y aplicaciones}

\subsection{Cálculo de autovalores y autovectores con NumPy}
\begin{lstlisting}[language=Python]
import numpy as np

A = np.array([[4, 1],
              [2, 3]])

autovalores, autovectores = np.linalg.eig(A)
print("Autovalores:", autovalores)
print("Autovectores:\n", autovectores)
# Nota: las columnas de autovectores son los vectores propios
# Verificar
for i in range(len(autovalores)):
    v = autovectores[:, i]
    lambda_i = autovalores[i]
    print(f"A @ v_{i} =", A @ v)
    print(f"lambda_{i} * v_{i} =", lambda_i * v)
    print()
\end{lstlisting}

\subsection{PCA desde cero en el dataset Iris}
\begin{lstlisting}[language=Python]
import numpy as np
import matplotlib.pyplot as plt
from sklearn.datasets import load_iris
from sklearn.preprocessing import StandardScaler

# Cargar datos
iris = load_iris()
X = iris.data
y = iris.target

# Centrar (PCA con matriz de covarianza, no necesita escalado si usamos correlación, pero por simplicidad centramos)
scaler = StandardScaler(with_std=False)  # solo centrar
X_centered = scaler.fit_transform(X)

# Calcular matriz de covarianza
cov = np.cov(X_centered.T)

# Autovalores y autovectores
eigvals, eigvecs = np.linalg.eig(cov)
# Ordenar de mayor a menor
idx = np.argsort(eigvals)[::-1]
eigvals = eigvals[idx]
eigvecs = eigvecs[:, idx]

print("Autovalores:", eigvals)
print("Proporción de varianza:", eigvals / np.sum(eigvals))

# Proyectar a 2 dimensiones
k = 2
W = eigvecs[:, :k]
X_pca = X_centered @ W

# Visualizar
plt.figure(figsize=(8,6))
for i in np.unique(y):
    plt.scatter(X_pca[y==i, 0], X_pca[y==i, 1], label=iris.target_names[i])
plt.xlabel('Componente 1')
plt.ylabel('Componente 2')
plt.legend()
plt.title('PCA de Iris (2 componentes)')
plt.grid(True)
plt.show()
\end{lstlisting}

\subsection{PCA con SVD (más estable numéricamente)}
Alternativamente, podemos usar SVD directamente, que es el método recomendado:

\begin{lstlisting}[language=Python]
U, S, Vt = np.linalg.svd(X_centered, full_matrices=False)
# Las componentes principales son las filas de Vt
# La proyección es X_centered @ Vt.T
X_pca_svd = X_centered @ Vt.T[:, :2]
# Los valores singulares al cuadrado / (n-1) son los autovalores de la covarianza
print("Autovalores desde SVD:", S**2 / (len(X)-1))
\end{lstlisting}

\subsection{PageRank simplificado}
Implementamos el método de la potencia para una pequeña matriz de enlaces.

\begin{lstlisting}[language=Python]
def power_iteration(A, num_simulations=100):
    # Estimación inicial
    v = np.random.rand(A.shape[1])
    v = v / np.linalg.norm(v)
    for _ in range(num_simulations):
        v = A @ v
        v = v / np.linalg.norm(v)
    # El autovalor asociado es (v^T A v) / (v^T v)
    eigenvalue = v.T @ A @ v
    return eigenvalue, v

# Ejemplo: matriz de transición (simétrica para simplificar)
A = np.array([[0.1, 0.2, 0.7],
              [0.3, 0.5, 0.2],
              [0.6, 0.3, 0.1]])
# Asegurar que las columnas suman 1 (estocástica)
A = A / A.sum(axis=0, keepdims=True)
print("Matriz estocástica:\n", A)

eigval, eigvec = power_iteration(A)
print("Autovalor principal (debe ser ~1):", eigval)
print("Autovector principal (ranking):", eigvec)
# Normalizar a suma 1 para que sea distribución
ranking = eigvec / eigvec.sum()
print("Distribución de importancia:", ranking)
\end{lstlisting}

\section{Notación en papers científicos}

En artículos de machine learning, es común encontrar expresiones como:

\begin{itemize}
    \item Descomposición espectral de una matriz simétrica: $\mathbf{A} = \mathbf{Q} \mathbf{\Lambda} \mathbf{Q}^\top$, con $\mathbf{Q}$ ortogonal.
    \item En PCA, la proyección se escribe $\mathbf{Z} = \mathbf{X} \mathbf{U}_k$, donde $\mathbf{U}_k$ contiene los $k$ autovectores principales.
    \item El teorema espectral para matrices de covarianza: $\boldsymbol{\Sigma} = \sum_{i=1}^d \lambda_i \mathbf{u}_i \mathbf{u}_i^\top$.
    \item En aprendizaje profundo, se estudia el espectro del Hessiano: $\nabla^2 \mathcal{L}(\mathbf{w})$, cuyos autovalores determinan la curvatura local.
    \item PageRank: $\mathbf{r} = \mathbf{G} \mathbf{r}$, con $\mathbf{G}$ matriz de Google.
\end{itemize}

Ejemplo típico:
\[
\mathbf{X}^\top \mathbf{X} \mathbf{v}_j = \lambda_j \mathbf{v}_j, \quad j=1,\dots,d
\]

\section{Ejercicios propuestos}

\begin{enumerate}
    \item Encuentra los autovalores y autovectores de $\mathbf{A} = \begin{pmatrix} 2 & 0 \\ 0 & -3 \end{pmatrix}$ e interpreta geométricamente.
    \item Demuestra que si $\mathbf{A}$ es simétrica, sus autovalores son reales. (Indicación: considera $\mathbf{v}^* \mathbf{A} \mathbf{v}$.)
    \item Aplica PCA al dataset de dígitos de sklearn (\texttt{load\_digits}) y visualiza las dos primeras componentes. ¿Se separan bien las clases?
    \item Investiga qué es el ``método de la potencia'' y por qué converge al autovector dominante. Implementa una versión con tolerancia.
    \item Lee la sección de análisis de autovalores en el paper original de PageRank (Brin \& Page, 1998) e identifica la notación.
\end{enumerate}

\section{Resumen}

En esta sesión hemos aprendido:
\begin{itemize}
    \item La definición de autovalor y autovector y su interpretación geométrica.
    \item Cómo calcularlos manualmente (caso 2x2) y con NumPy.
    \item Propiedades fundamentales (traza, determinante, diagonalización).
    \item Aplicaciones estrella: PCA (reducción de dimensionalidad), PageRank (ranking de páginas) y análisis de redes neuronales.
    \item Implementación práctica en Python con datasets reales.
    \item La notación habitual en la literatura científica.
\end{itemize}
Estos conceptos son la base de muchas técnicas modernas de machine learning y análisis de datos.

\end{document}